\chapter{Introducción}
La fiabilidad de un sistema de información depende fundamentalmente de la solidez de su base de datos frente a errores humanos o lógicos. Para la comercializadora Impor Tech S.A.S., con su alto volumen de transacciones diarias en ventas y control de inventarios se debe realizar varias pruebas técnicas y escenarios de estrés y fallos intencionales. El presente documento expondrá las pruebas de validación aplicadas al esquema relacional de la compañía para asegurar la calidad de la información.

\section{Objetivos del Trabajo}
\begin{itemize}
    \item \textbf{Objetivo General}: Validar el diseño y el correcto funcionamiento del modelo de base de datos de Import Tech S.A.S. mediante la ejecución de pruebas de integridad de datos y restricciones relacionales.
    \item \textbf{Objetivo Específicos}:
        \begin{enumerate}
            \item Ejecutar sentencias SQL que reten las restricciones de integridad de entidad (\colorbox{gray!20}{\textcolor{lineText}{\texttt{PRIMARY KEY}}}) e integridad de dominio (\colorbox{gray!20}{\textcolor{lineText}{\texttt{NOT NULL}}}).   
            \item Comprobar la integridad referencial intentando vulnerar las relaciones de llaves foráneas (\colorbox{gray!20}{\textcolor{lineText}{\texttt{FOREIGN KEY}}}) entre entidades dependientes e independientes.
            \item Documentar la respuesta de la base de datos para certificar la consistencia transaccional del sistema.
        \end{enumerate}
\end{itemize}